\akapit Badania przeprowadzono w tunelu aerodynamicznym wyposażonym w wagę aerodynamiczną Witoszyńskiego (typu JAW). Badany model zawieszony jest pionowo na napiętych drutach wewnątrz obrotowej ramy, której kąt natarcia można regulować w zakresie ±30° za pomocą samohamownej przekładni ślimakowej (odczyt kąta odbywa się ze skali na rurze wagi lub za pomocą przetwornika potencjometrycznego). Konstrukcja stanowiska umożliwia oddzielny (sekwencyjny) pomiar trzech składowych aerodynamicznych: siły nośnej, siły oporu oraz momentu pochylającego $M_yo$. Wybór mierzonej wielkości realizowany jest poprzez zablokowanie odpowiedniej pary sworzni $(S_z; S_x; S_m)$ i zwolnienie pozostałych, co nadaje ramie swobodę obrotu względem wybranej z trzech wzajemnie prostopadłych osi. Obciążenia z modelu przekazywane są przez układ drutów, ramę i dźwignie dwuramienne na tensometryczny przetwornik siły oraz połączoną z nim równolegle wagę szalkową służącą do wzorcowania układu. 

\akapit Badania przeprowadzono w tunelu aerodynamicznym, choć stanowisko to jest bazowo wyposażone w klasyczną wagę aerodynamiczną Witoszyńskiego (typu JAW), w trakcie realizowanego ćwiczenia zastosowano zmodyfikowany, zautomatyzowany układ mocowania i pomiaru. Zamiast tradycyjnego podwieszenia modelu na drutach, do obrotowej ramy wagi wykorzystano tzw. żądło (odpowiednio wygięty wysięgnik prętowy), na którego końcu osadzono badany model samolotu odrzutowego. Do wyznaczenia obciążeń aerodynamicznych wykorzystano, wewnętrzną wagę tensometryczną umieszczoną bezpośrednio wewnątrz modelu, co wyeliminowało konieczność stosowania zewnętrznego układu dźwigni i drutów. Zmiana kąta natarcia (realizowana poprzez obrót ramy z żądłem za pomocą przekładni napędzanej silnikiem elektrycznym) oraz akwizycja danych pomiarowych z czujnika odbywały się zdalnie z poziomu komputera wyposażonego w dedykowane oprogramowanie pomiarowe. 
